% ==========================================================================
%  Prova 2 (15/07/2026) — Gabarito Comentado
%  Macroeconomia I (PPGEco-UFES, 2026)
%  Prof. Dr. Ricardo Ramalhete Moreira
%  Cobre: Consumo (cap.8), Investimento (cap.9),
%         Inflação e Política Monetária (cap.11), Política Fiscal (cap.12)
%  Prova transcrita das fotos autorizadas; gabarito reconstruído a partir
%  das chaves oficiais (P2 2024, P2 2021, Lista 3 2026) e slides — NÃO é
%  o gabarito oficial do docente. Q9 e Q12 têm nota de ambiguidade.
% ==========================================================================
\documentclass[12pt, a4paper]{article}

\usepackage[utf8]{inputenc}
\usepackage[T1]{fontenc}
\usepackage{lmodern}
\usepackage[brazil]{babel}
\usepackage{microtype}

\usepackage[top=2.5cm, bottom=2.5cm, left=3cm, right=3cm, headheight=15pt]{geometry}

\usepackage{amsmath, amssymb, amsthm}
\usepackage{mathtools}
\usepackage{bm}

\usepackage{booktabs}
\usepackage{array}
\usepackage{multicol}
\usepackage{xcolor}
\usepackage{hyperref}
\usepackage{enumitem}
\usepackage{tcolorbox}
\tcbuselibrary{breakable}
\usepackage{fancyhdr}
\usepackage{titlesec}

% --------------------------------------------------------------------------
% Cores
% --------------------------------------------------------------------------
\definecolor{cyanrbc}{HTML}{3b9dbd}
\definecolor{orangerbc}{HTML}{d97b39}
\definecolor{grayrbc}{HTML}{6b7280}
\definecolor{greenbox}{HTML}{166534}
\definecolor{greenbg}{HTML}{dcfce7}
\definecolor{boxbg}{HTML}{f5f0e8}
\definecolor{boxborder}{HTML}{3b9dbd}
\definecolor{provabg}{HTML}{fff7ed}
\definecolor{provaborder}{HTML}{d97b39}
\definecolor{correctbg}{HTML}{dcfce7}
\definecolor{correctborder}{HTML}{166534}

% --------------------------------------------------------------------------
% Caixas
% --------------------------------------------------------------------------
\newtcolorbox{resultbox}[1][]{standard, colback=boxbg, colframe=boxborder,
  boxrule=1pt, arc=3pt, left=6pt, right=6pt, top=4pt, bottom=4pt,
  fonttitle=\bfseries\small, breakable, #1}
\newtcolorbox{provabox}[1][]{standard, colback=provabg, colframe=provaborder,
  boxrule=1.2pt, arc=3pt, left=6pt, right=6pt, top=4pt, bottom=4pt,
  fonttitle=\bfseries\small, breakable, #1}
\newtcolorbox{correctbox}[1][]{standard, colback=correctbg, colframe=correctborder,
  boxrule=1pt, arc=3pt, left=6pt, right=6pt, top=4pt, bottom=4pt,
  fonttitle=\bfseries\small\color{correctborder}, breakable, #1}

% --------------------------------------------------------------------------
% Comandos matemáticos
% --------------------------------------------------------------------------
\newcommand{\E}{\mathbb{E}}
\newcommand{\dif}{\mathrm{d}}
\newcommand{\VAR}{\mathrm{VAR}}

\setlength{\parskip}{0.5\baselineskip}
\setlength{\parindent}{0pt}

\titleformat{\section}{\large\bfseries\color{cyanrbc}}{Questão \thesection.}{0.5em}{}[\vspace{-0.3em}\rule{\linewidth}{0.4pt}]
\titleformat{\subsection}{\normalsize\bfseries\color{grayrbc}}{}{0em}{}

\pagestyle{fancy}
\fancyhf{}
\lhead{\footnotesize\textcolor{grayrbc}{P2 15/07/2026 --- Gabarito Comentado --- Macroeconomia I (PPGEco-UFES)}}
\rhead{\footnotesize\textcolor{grayrbc}{Consumo $\cdot$ Investimento $\cdot$ Moeda $\cdot$ Fiscal}}
\cfoot{\small\thepage}
\renewcommand{\headrulewidth}{0.4pt}

% ==========================================================================
\begin{document}
% ==========================================================================

% --------------------------------------------------------------------------
% Capa
% --------------------------------------------------------------------------
\begin{titlepage}
\begin{center}
\vspace*{1.2cm}
{\Large\bfseries Universidade Federal do Espírito Santo}\\[0.5em]
{\large Programa de Pós-Graduação em Economia (PPGEco-UFES)}\\[2em]

\rule{\linewidth}{1.2pt}\\[0.8em]
{\LARGE\bfseries\color{cyanrbc} Prova 2 --- Gabarito Comentado\\[0.4em]
\large Macroeconomia I --- 15/07/2026}\\[0.4em]
{\normalsize\color{grayrbc} Romer caps.\ 8, 9, 11 e 12}\\[0.8em]
\rule{\linewidth}{1.2pt}

\vspace{1em}
{\large Prof.\ Dr.\ Ricardo Ramalhete Moreira}

\vspace{1.5em}
\begin{tcolorbox}[standard, colback=provabg, colframe=provaborder, width=0.88\textwidth,
  arc=4pt, boxrule=1.2pt, left=8pt, right=8pt, top=6pt, bottom=6pt]
\centering
\textbf{\textcolor{provaborder}{Conteúdo:}} 12 questões objetivas (escolha única)\\[0.4em]
\textbf{Consumo} (Q1--Q3) $\cdot$ \textbf{Investimento / $q$ de Tobin} (Q4--Q6)\\[0.2em]
\textbf{Política Fiscal} (Q7--Q8, Q10) $\cdot$ \textbf{Síntese} (Q9) $\cdot$
\textbf{Política Monetária} (Q11--Q12)
\end{tcolorbox}

\vspace{1.5em}
\begin{tcolorbox}[standard, colback=correctbg, colframe=correctborder,
  width=0.88\textwidth, arc=4pt, boxrule=1pt, left=8pt, right=8pt,
  top=6pt, bottom=6pt]
\centering
\textbf{\textcolor{correctborder}{Gabarito Rápido (reconstruído)}}\\[0.6em]
\begin{tabular}{cccccc}
\toprule
Q1 & Q2 & Q3 & Q4 & Q5 & Q6 \\
\textbf{ii} & \textbf{iii} & \textbf{iii} & \textbf{iv} & \textbf{iv} & \textbf{i} \\
\midrule
Q7 & Q8 & Q9 & Q10 & Q11 & Q12 \\
\textbf{iii} & \textbf{ii} & \textbf{iv}$^{*}$ & \textbf{iv} & \textbf{ii} & \textbf{ii}$^{*}$ \\
\bottomrule
\end{tabular}\\[0.5em]
{\footnotesize $^{*}$ Q9 e Q12 têm nota de ambiguidade de redação --- ver comentários.}
\end{tcolorbox}

\vspace{1.2em}
\begin{tcolorbox}[standard, colback=boxbg, colframe=boxborder, width=0.88\textwidth,
  arc=4pt, boxrule=1pt, left=8pt, right=8pt, top=5pt, bottom=5pt]
\raggedright\small
\textbf{Mapa de reciclagem} (de onde cada questão veio):
Q2 = P2 2024 Q2 $\cdot$ Q3 = P2 2021/2024 Q1 \emph{com itens invertidos}
$\cdot$ Q4 = P2 2021 Q2 / P2 2024 Q3 $\cdot$ Q5 = diagrama de fases (Lista 3
Q5) $\cdot$ Q6 = P2 2024 Q4 $\cdot$ Q7 = P2 2024 Q6 $\cdot$ Q8 = P2 2024 Q5
$\cdot$ Q11 = Lista 3 2026 Q6 (Clarida) $\cdot$ Q12 = P2 2021 Q3 reformulada
(estrutura a termo) $\cdot$ Q9 e Q10 = inéditas (síntese e slides da Aula 13).
\end{tcolorbox}

\vfill
{\small Gabarito \textbf{reconstruído} a partir das chaves oficiais da P2 2024
e da Lista 3 2026, das respostas da P2 2021 e dos slides das Aulas 10--13.
Não substitui o gabarito oficial do docente.}
\end{center}
\end{titlepage}

% ==========================================================================
\begin{center}
{\large\bfseries\color{cyanrbc} Bloco I --- Teoria do Consumo (Cap.\ 8)}\\[0.2em]
\rule{\linewidth}{0.6pt}
\end{center}

\setcounter{section}{0}

% ------------------------------------------------------------------
\section{Equação de Euler: Juros e Crescimento do Consumo}
% ------------------------------------------------------------------

\textbf{Enunciado.} Considere a condição de Euler proveniente de um modelo
microfundamentado de consumo intertemporal:
\[
\frac{C_{t+1}}{C_t} = \left(\frac{1+r}{1+\rho}\right)^{\frac{1}{\theta}}.
\]
Com base no modelo e na evidência empírica da literatura, analise as
afirmativas abaixo:

\begin{enumerate}[label=\Alph*.]
\item Quanto maior o parâmetro $\theta$, menor tende a ser a resposta do
crescimento do consumo às variações da taxa real de juros.
\item As evidências empíricas de Hansen e Singleton (1983), Hall (1988b) e
Campbell e Mankiw (1989) sugerem que o crescimento do consumo apresenta baixa
sensibilidade às variações da taxa real de juros.
\item Mantidos constantes os demais parâmetros, um aumento da taxa real de
juros eleva a taxa de crescimento do consumo ao longo do tempo, refletindo o
efeito substituição intertemporal.
\item A taxa de desconto subjetiva ($\rho$) decorre diretamente da restrição
orçamentária intertemporal do consumidor, sendo determinada pelas condições de
mercado.
\end{enumerate}
Alternativas: i) apenas A e B; ii) apenas A, B e C; iii) apenas B, C e D;
iv) todas.

\begin{correctbox}[title={Resposta: ii) Apenas A, B e C estão corretas}]
\textbf{A --- V.} A sensibilidade do crescimento do consumo a $r$ é a
elasticidade de substituição intertemporal, $1/\theta$: $\theta$ maior
$\Rightarrow$ resposta menor.

\textbf{B --- V.} É o fato estilizado da literatura (chave da Lista 3, Q2):
baixo poder explicativo dos juros reais, atribuído a $\theta$ elevado.

\textbf{C --- V.} $r\uparrow$ com $\rho$ e $\theta$ dados $\Rightarrow$
$(1+r)/(1+\rho)$ maior $\Rightarrow$ trajetória de consumo mais inclinada ---
poupar fica mais atrativo (efeito substituição intertemporal).

\textbf{D --- F.} $\rho$ é um \emph{parâmetro de preferências} (impaciência),
primitivo do modelo; não ``decorre'' da restrição orçamentária nem é
``determinado pelas condições de mercado''. Quem vem do mercado é $r$.
\end{correctbox}

% ------------------------------------------------------------------
\section{Otimização do Consumo com Governo}
% ------------------------------------------------------------------

\textbf{Enunciado.} Seja um modelo de otimização do consumo com a taxa real de
juros, taxa de desconto subjetiva e aversão ao risco, além de uma restrição
orçamentária intertemporal padrão. Julgue os itens, escolhendo o correto:

\begin{enumerate}[label=\roman*)]
\item Como é de se esperar, a variação do consumo no tempo é inversamente
relacionada à taxa real de juros;
\item O fato estilizado reportado pela literatura é claro: há grande impacto
das taxas reais de juros no comportamento do consumo das famílias;
\item A restrição orçamentária das famílias, com presença do governo, não
sofre influência das proporções de T e D, sendo os níveis reais T para
arrecadação e D para dívida do Governo;
\item Tanto a elevação da taxa de desconto, quanto do grau de aversão ao
risco, levam a maiores graus de eficácia de juros reais sobre a mudança do
consumo.
\end{enumerate}

\begin{correctbox}[title={Resposta: iii) --- reciclada: gabarito oficial da P2 2024 (Q2-iv)}]
\textbf{iii --- V.} É a equivalência ricardiana na restrição das famílias:
substituindo a restrição do governo (satisfeita com igualdade), $T$ e $D$
desaparecem e resta apenas o valor presente de $G$:
\[
\int_0^\infty e^{-R(t)}C\,\dif t \leq K(0) + \int_0^\infty e^{-R(t)}W\,\dif t
- \int_0^\infty e^{-R(t)}G\,\dif t.
\]

\textbf{Por que as demais são falsas:} i) a relação é \emph{positiva} (Euler);
ii) o fato estilizado é de impacto \emph{fraco}; iv) maior $\theta$
\emph{reduz} a sensibilidade ($1/\theta$) --- e $\rho$ não a eleva.
\end{correctbox}

% ------------------------------------------------------------------
\section{Renda Permanente: Campbell-Mankiw e Extensões}
% ------------------------------------------------------------------

\textbf{Enunciado.} Na teoria do consumo, existem aspectos que não podem ser
explicados tão somente com base na hipótese da renda permanente. Julgue os
itens, escolhendo o correto:

\begin{enumerate}[label=\roman*)]
\item Seja a regressão proposta em Campbell e Mankiw (1989):
$C_t - C_{t-1} = \lambda(Y_t - Y_{t-1}) + (1-\lambda)e_t$. A rejeição plena da
hipótese do passeio aleatório requer encontrar valores convergentes com
$\lambda = 0$;
\item A teoria da renda permanente pode ser entendida como um caso geral,
sendo a previsão keynesiana do consumo um caso particular. Em amostras longas,
a variância da renda transitória domina a variância da renda permanente. Isto
explica o menor valor da propensão a consumir no longo prazo;
\item A correlação entre taxas reais de juros e crescimento do consumo, em
modelos de maximização das famílias, não implica o abandono da hipótese da
renda permanente;
\item A poupança por motivo precaucional está associada à incerteza quanto ao
futuro, tal como inferido em
$\E[g^c] \simeq \tfrac{1}{2}(\theta+1)\VAR(g^c)$. [\ldots] Isto ajuda a
entender o fato estilizado de uma elevada poupança pelas famílias.
\end{enumerate}

\begin{correctbox}[title={Resposta: iii) --- P2 2021/2024 Q1 com \emph{dupla inversão} de itens}]
O professor pegou a questão reciclada de 2021/2024 e \textbf{inverteu os dois
itens decisivos}:

\textbf{i --- F (era V com $\lambda = 1$).} A rejeição \emph{plena} do passeio
aleatório exige $\lambda = 1$ (todas as famílias na regra de bolso). Com
$\lambda = 0$ a hipótese de Hall é \emph{confirmada}, não rejeitada. Quem
decorou a letra da resposta antiga sem entender caiu aqui.

\textbf{iii --- V (era F sem o ``não'').} A equação de Euler com juros é uma
\emph{extensão} do mesmo arcabouço otimizador da renda permanente --- a
correlação juros--consumo não exige abandonar a hipótese.

\textbf{Por que as demais são falsas:} ii) em amostras longas domina a
variância \emph{permanente} (e a propensão estimada aproxima-se de 1 ---
$\hat b = \VAR(Y^P)/[\VAR(Y^P)+\VAR(Y^T)]$); iv) o fato estilizado é de
poupança \emph{baixa} (puzzle resolvido por Carroll com a alta taxa de
desconto).
\end{correctbox}

% ==========================================================================
\begin{center}
{\large\bfseries\color{cyanrbc} Bloco II --- Investimento e $q$ de Tobin (Cap.\ 9)}\\[0.2em]
\rule{\linewidth}{0.6pt}
\end{center}

% ------------------------------------------------------------------
\section{Teoria do Investimento}
% ------------------------------------------------------------------

\textbf{Enunciado.} Quanto aos aspectos ligados à teoria do investimento,
julgue qual a verdadeira:

\begin{enumerate}[label=\roman*)]
\item Considerando um modelo com custos de ajuste do estoque de capital, a
condição de maximização de lucros implica que a firma investe até o ponto em
que o valor do capital supera o custo de sua aquisição;
\item A partir do diagrama de fase para o $q$ de Tobin e o estoque de capital
$k$, o \emph{locus} abaixo de $\Delta K = 0$ e à direita de $\Delta q = 0$ é
uma região de divergência crescente em relação ao equilíbrio de longo prazo da
indústria;
\item Enquanto uma redução permanente do produto vem acompanhada de diminuição
do investimento, dada a queda do $q$ de Tobin, uma redução transitória do
produto possui efeito nulo sobre o investimento;
\item Pela equação de variação do $q$ de Tobin, uma situação de estoque de
capital da indústria muito baixo, acompanhada de nível também muito reduzido
para o $q$, implica uma queda adicional do último.
\end{enumerate}

\begin{correctbox}[title={Resposta: iv) --- reciclada: gabarito oficial P2 2021 (Q2) e P2 2024 (Q3)}]
\textbf{iv --- V.} Por $\dot q = rq - \pi(K)$: $q$ muito baixo torna $rq$
pequeno; $K$ muito baixo torna $\pi(K)$ grande $\Rightarrow \dot q < 0$: o $q$
cai ainda mais (região divergente a sudoeste da trajetória de sela).

\textbf{Por que as demais são falsas:} i) no ótimo o valor \emph{iguala} o
custo ($1 + C'(I) = q$), nunca ``supera''; ii) essa região contém o braço
\emph{convergente} da sela vindo do sudeste; iii) o efeito do choque
transitório é \emph{menor}, não nulo.
\end{correctbox}

% ------------------------------------------------------------------
\section{Diagrama de Fases: Sinais dos Quadrantes}
% ------------------------------------------------------------------

\textbf{Enunciado.} Ainda com respeito ao modelo do $q$ de Tobin, julgue qual
a verdadeira:

\begin{enumerate}[label=\roman*)]
\item Na análise de diagrama de fase, com base nos eixos $\Delta q = 0$ /
$\Delta K = 0$, uma situação acima de $\Delta K = 0$ e à esquerda de
$\Delta q = 0$ leva a divergência frente ao equilíbrio da indústria;
\item uma situação abaixo de $\Delta K = 0$ e à esquerda de $\Delta q = 0$
leva a $\Delta q > 0$ e $\Delta K < 0$;
\item uma situação abaixo de $\Delta K = 0$ e à direita de $\Delta q = 0$
leva a $\Delta q < 0$ e $\Delta K < 0$;
\item uma situação acima de $\Delta K = 0$ e à direita de $\Delta q = 0$
leva a $\Delta q > 0$ e $\Delta K > 0$.
\end{enumerate}

\begin{correctbox}[title={Resposta: iv) --- puro sinal dos quadrantes}]
Regras: acima de $\Delta K = 0$ $\Leftrightarrow q > 1 \Rightarrow
\Delta K > 0$; abaixo $\Rightarrow \Delta K < 0$. À direita de
$\Delta q = 0$ $\Leftrightarrow rq > \pi(K) \Rightarrow \Delta q > 0$; à
esquerda $\Rightarrow \Delta q < 0$.

\textbf{iv --- V.} Acima + direita (quadrante NE): $\Delta K > 0$ e
$\Delta q > 0$. \checkmark\ (Região divergente: tudo cresce.)

\textbf{Por que as demais são falsas:} i) acima + esquerda é o quadrante NO
--- região da \emph{trajetória de sela} (chave da Lista 3, Q5): converge a
$E$, não diverge; ii) à esquerda $\Delta q < 0$, não $> 0$; iii) à direita
$\Delta q > 0$, não $< 0$.
\end{correctbox}

% ------------------------------------------------------------------
\section{$q$ de Tobin: Juros Longos}
% ------------------------------------------------------------------

\textbf{Enunciado.} Com respeito ao modelo do $q$ de Tobin, julgue qual a
verdadeira:

\begin{enumerate}[label=\roman*)]
\item \emph{Coeteris paribus}, mais importante do que os juros reais no
momento $t$, é a expectativa de $r$ para o longo prazo, no que respeita ao $q$
da empresa;
\item A expressão de valor da empresa, tal como representada por
\[
q(t) = \int_{0}^{\infty} e^{-rt}\,\pi(K(t))\,\dif t,
\]
implica que o valor de mercado das empresas está positivamente relacionado ao
estoque de capital da indústria;
\item Uma expectativa de mercado segundo a qual a taxa de inflação deve
aumentar frente à meta inflacionária nos próximos anos é condizente com maior
$q$ da empresa, \emph{coeteris paribus}, já que isto é acompanhado de
expectativa de maior crescimento do PIB e, desta forma, de maiores fluxos
futuros de caixa para as empresas;
\item À medida que o estoque de capital da indústria $K$ aumenta, maior é o
valor $q$, de modo que a função $\Delta q = 0$ é positivamente inclinada no
eixo $q \times K$.
\end{enumerate}

\begin{correctbox}[title={Resposta: i) --- reciclada: gabarito oficial da P2 2024 (Q4-iii)}]
\textbf{i --- V.} O $q$ é o valor presente de \emph{todo} o fluxo futuro de
lucros: a taxa de desconto relevante é a estrutura de juros esperada
(sobretudo os juros reais \emph{longos}), não a taxa instantânea de hoje.

\textbf{Por que as demais são falsas:} ii) $\pi'(K) < 0$: relação
\emph{negativa}; iii) inflação esperada acima da meta $\Rightarrow$ Banco
Central eleva o juro \emph{real} ($\phi_\pi > 1$) $\Rightarrow$ desconto maior
$\Rightarrow$ $q$ \emph{menor}; iv) como $\pi'(K)<0$, o \emph{locus}
$\Delta q = 0$ é \emph{negativamente} inclinado.
\end{correctbox}

% ==========================================================================
\begin{center}
{\large\bfseries\color{cyanrbc} Bloco III --- Política Fiscal e Síntese (Cap.\ 12)}\\[0.2em]
\rule{\linewidth}{0.6pt}
\end{center}

% ------------------------------------------------------------------
\section{Restrição Orçamentária do Governo e Solvência}
% ------------------------------------------------------------------

\textbf{Enunciado.} A respeito da restrição orçamentária do governo e de suas
implicações em termos fiscais e econômicos, julgue qual a verdadeira:

\begin{enumerate}[label=\roman*)]
\item A avaliação quanto à sustentabilidade da dívida pública (sua solvência)
é feita a partir dos resultados fiscais correntes;
\item Dois países A e B possuem uma dívida pública real de mesma magnitude
(medida em uma moeda comum); no entanto, os saldos fiscais correntes de A são
muito superiores aos de B ($T_{tA} - G_{tA} > T_{tB} - G_{tB}$). Logo, a
dívida pública de B será necessariamente considerada mais arriscada que a de
A;
\item Os superávits primários corrente e futuros esperados precisam cobrir a
dívida pública atual em termos reais a fim de que a restrição orçamentária
seja satisfeita;
\item A mensuração e avaliação quanto à adequação dos saldos fiscais não sofre
qualquer influência da dinâmica inflacionária, visto que seus efeitos são
anulados necessariamente pela política monetária.
\end{enumerate}

\begin{correctbox}[title={Resposta: iii) --- reciclada: gabarito oficial da P2 2024 (Q6-iii)}]
\textbf{iii --- V.} É a condição de solvência:
$\int_0^\infty e^{-R(t)}[T(t)-G(t)]\,\dif t \geq D(0)$ --- superávits
primários corrente \emph{e futuros esperados}, a valor presente, cobrindo a
dívida atual.

\textbf{Por que as demais são falsas:} i) solvência se avalia pelos fluxos
\emph{projetados} (o corrente é fração mínima do somatório); ii) o
``necessariamente'' derruba --- expectativas de fluxos futuros podem tornar B
menos arriscado; iv) dupla generalização indevida (``não sofre qualquer'' /
``anulados necessariamente'').
\end{correctbox}

% ------------------------------------------------------------------
\section{Política Fiscal: Equivalência Ricardiana}
% ------------------------------------------------------------------

\textbf{Enunciado.} Considere as seguintes afirmativas acerca da política
fiscal, escolhendo apenas a verdadeira:

\begin{enumerate}[label=\roman*)]
\item A condição de restrição orçamentária intertemporal do governo estabelece
que o somatório das receitas fiscais (líquidas das transferências), a valor
presente, precisa ser maior que a dívida pública real no momento atual, mais o
fluxo esperado das despesas do governo, a valor presente;
\item A variação do endividamento público é entendida como resultado dos
déficits fiscais. Neste sentido, uma vez que o governo satisfaça a restrição
de orçamento no tempo, a substituição de financiamento via tributos por mais
dívida não possui qualquer impacto sobre o comportamento do consumidor;
\item Uma política de redução de alíquotas tributárias, ao elevar a renda
disponível corrente das famílias, vem acompanhada de aumento do consumo e
desta maneira de um estímulo real positivo para o ciclo econômico;
\item Os gastos públicos corrente e futuros esperados não têm impacto direto
na restrição orçamentária das famílias, embora impactem o equilíbrio dinâmico
de longo prazo da economia como um todo.
\end{enumerate}

\begin{correctbox}[title={Resposta: ii) --- reciclada: gabarito oficial da P2 2024 (Q5-iii)}]
\textbf{ii --- V.} Equivalência ricardiana: dada a satisfação da restrição do
governo, o mix tributos/dívida não entra na restrição das famílias nem nas
preferências --- logo não altera o consumo. (Este é o caso raro em que ``não
possui qualquer impacto'' é \emph{verdadeiro} --- a exceção à regra de bolso
de desconfiar de afirmações absolutas.)

\textbf{Por que as demais são falsas:} i) a condição é de ``maior \emph{ou
igual}'' ($\geq$) --- a satisfação \emph{com igualdade} é admissível (e é a
usada na derivação da equivalência); o ``precisa ser maior'' estrito foi
considerado falso na chave de 2024; iii) sob equivalência ricardiana a renda
extra do corte é \emph{poupada}; iv) $G$ entra \emph{diretamente} na restrição
das famílias após a substituição ($-\int e^{-R}G$).
\end{correctbox}

% ------------------------------------------------------------------
\section{Síntese da Disciplina}
% ------------------------------------------------------------------

\textbf{Enunciado.} Do exposto ao longo da disciplina, deduz-se que:

\begin{enumerate}[label=\roman*)]
\item os níveis de investimento e de renda per capita a longo prazo podem ser
estimulados positivamente a partir de expansão do consumo público;
\item existem restrições ao aumento de renda (e consumo) per capita a longo
prazo, dentre os quais aqueles ligados ao comportamento das famílias. Exemplo
é a aversão ao risco e a taxa de desconto. Ambos relacionam-se positivamente
com a renda per capita;
\item a elevação dos investimentos privados, essenciais no planejamento de
crescimento e bem-estar de um país, requer a viabilidade econômico-financeira
daqueles primeiros. Neste sentido, a elevação do produto, ainda que
transitória, leva os setores produtivos a níveis mais elevados de capital no
longo prazo, melhorando persistentemente as condições de vida da população;
\item é possível sugerir que uma redução da taxa real de juros de curto prazo
seja acompanhada de queda dos investimentos, caso, ao mesmo tempo, haja piora
das projeções de inflação futura, dado que os juros reais de longo prazo
tendem a aumentar em tal cenário.
\end{enumerate}

\begin{correctbox}[title={Resposta: iv) --- com nota de ambiguidade na redação do item ii}]
\textbf{iv --- V.} É a cadeia completa do curso: o $q$ (e o investimento)
responde aos juros reais \emph{longos}, que são média das taxas curtas
esperadas. Piora das projeções de inflação + Princípio de Taylor
($\phi_\pi>1$) $\Rightarrow$ mercado espera juros reais futuros maiores
$\Rightarrow$ juros longos sobem $\Rightarrow$ $q$ cai $\Rightarrow$
investimento cai --- mesmo com a taxa curta corrente em queda.

\textbf{Por que as demais são falsas:} i) expansão do consumo público não
estimula investimento/renda de longo prazo (em RCK desloca $c$, não $k^*$;
com equivalência ricardiana, $G$ reduz os recursos das famílias);
iii) elevação \emph{transitória} do produto tem efeito temporário e
\emph{menor} --- não eleva o capital de longo prazo ``persistentemente''
(Fig.\ 9.6); \textbf{ii)} aversão ao risco e taxa de desconto relacionam-se
\emph{negativamente} com a renda per capita de longo prazo --- pela regra de
ouro modificada, $f'(k^*) = \rho + \theta g$: maiores $\theta$ ou $\rho$
$\Rightarrow$ $k^*$ e $y^*$ \emph{menores}. A primeira parte do item é
verdadeira (são restrições ligadas às famílias), mas o ``positivamente''
final o torna falso.

\emph{Nota de ambiguidade}: se a intenção do docente era ``negativamente''
(erro de redação), o item ii seria o correto. Como escrito, iv é a única
afirmativa integralmente verdadeira.
\end{correctbox}

% ------------------------------------------------------------------
\section{Ajustes Fiscais: Evidência Empírica}
% ------------------------------------------------------------------

\textbf{Enunciado.} Seja a literatura empírica sobre efeitos dos programas de
ajuste fiscal, quanto aos resultados empíricos a partir dos mesmos.
Identifique a afirmação que se adequa às implicações dessa literatura:

\begin{enumerate}[label=\roman*)]
\item Pode-se dizer que tal literatura sugere a invalidade da hipótese da
renda permanente;
\item Os resultados sugerem que ajustes fiscais são acompanhados por
multiplicadores fiscais positivos;
\item As expectativas dos agentes econômicos são consideradas como exógenas, e
o resultado da política fiscal sobre a atividade ampara-se na relação dos
dispêndios com a renda corrente;
\item Sugerem a possibilidade de multiplicadores fiscais negativos.
\end{enumerate}

\begin{correctbox}[title={Resposta: iv) --- contrações fiscais expansionistas (slides da Aula 13)}]
\textbf{iv --- V.} Giavazzi e Pagano (1990) --- Dinamarca e Irlanda nos anos
1980 --- e Alesina e Perotti (1997): consolidações fiscais podem vir
acompanhadas de \emph{expansão} da atividade (booms de consumo), via
expectativas (juros menores, menos distorções futuras, menor risco de crise).
Corte de gasto com produto \emph{subindo} $\Rightarrow$ multiplicador fiscal
\emph{negativo} --- exatamente ``a possibilidade'' que a literatura sugere.

\textbf{Por que as demais são falsas:} i) o mecanismo \emph{usa} a renda
permanente (renda futura esperada move consumo corrente), não a invalida;
ii) multiplicadores \emph{positivos} são o caso keynesiano padrão --- o
oposto da implicação distintiva dessa literatura; iii) as expectativas são o
\emph{centro endógeno} do mecanismo, não exógenas.
\end{correctbox}

% ==========================================================================
\begin{center}
{\large\bfseries\color{cyanrbc} Bloco IV --- Política Monetária (Cap.\ 11)}\\[0.2em]
\rule{\linewidth}{0.6pt}
\end{center}

% ------------------------------------------------------------------
\section{Regra de Taylor Forward-Looking e Credibilidade}
% ------------------------------------------------------------------

\textbf{Enunciado.} Considere uma economia na qual o Banco Central conduz a
política monetária sob um regime de metas de inflação. Suponha que a taxa
nominal de juros seja determinada segundo uma regra de Taylor simplificada,
\[
i_t = r^n + \E_t[\pi_{t+n}] + \phi_\pi\bigl(\E_t[\pi_{t+n}] - \pi^*\bigr)
+ \phi_y \E_t\bigl[\ln Y_{t+n} - \ln Y^n_{t+n}\bigr],
\qquad \phi_\pi > 1.
\]
Em determinado período ocorre um choque de demanda que eleva a inflação
esperada acima da meta, sem alterar inicialmente o hiato esperado do produto.
Analise as afirmativas abaixo:

\begin{enumerate}[label=\Alph*.]
\item Como $\phi_\pi > 1$, o Banco Central eleva a taxa nominal de juros mais
do que proporcionalmente ao aumento da inflação, provocando aumento da taxa
real de juros.
\item O aumento da taxa real de juros tende a reduzir a demanda agregada,
contribuindo para a convergência da inflação esperada à meta.
\item Sob elevada credibilidade, parte do ajuste ocorre por meio da redução
das expectativas de inflação futura, diminuindo o custo desinflacionário.
\item Quanto maior a credibilidade do Banco Central, maior tende a ser a
persistência inflacionária após o choque.
\end{enumerate}
Alternativas: i) apenas A e B; ii) apenas A, B e C; iii) apenas B e D;
iv) todas.

\begin{correctbox}[title={Resposta: ii) Apenas A, B e C --- Princípio de Taylor (chave da Lista 3, Q6)}]
\textbf{A --- V.} $\Delta r_t = \Delta i_t - \Delta\E_t(\pi_{t+n}) > 0$ quando
$\phi_\pi > 1$: o juro real sobe.

\textbf{B --- V.} Juro real maior contrai a demanda agregada e faz a inflação
convergir à meta (política contracíclica).

\textbf{C --- V.} Com credibilidade, as expectativas se ancoram e parte da
desinflação ocorre ``de graça'' (menor custo desinflacionário em produto).

\textbf{D --- F.} É o inverso: credibilidade \emph{reduz} a persistência
inflacionária --- $\phi_\pi < 1$ (ou credibilidade baixa) é que gera
persistência e desancoragem.
\end{correctbox}

% ------------------------------------------------------------------
\section{Fisher e Estrutura a Termo}
% ------------------------------------------------------------------

\textbf{Enunciado.} Considere a equação de Fisher, $i = r + \pi^e$, e a
hipótese de que a estrutura a termo das taxas de juros incorpora as
expectativas dos agentes sobre a trajetória futura da política monetária.
Analise as proposições:

\begin{enumerate}[label=\Alph*.]
\item Um aumento permanente da taxa de crescimento da oferta monetária tende a
elevar a inflação esperada e, no longo prazo, também as taxas nominais de
juros.
\item Caso o Banco Central anuncie de forma crível uma política monetária
permanentemente mais restritiva, as taxas de juros de longo prazo podem cair
imediatamente, mesmo antes da redução efetiva da inflação.
\item Se os agentes antecipam inflação futura menor, a queda das expectativas
inflacionárias pode reduzir as taxas nominais de longo prazo
independentemente da taxa Selic corrente.
\item A existência de expectativas racionais implica que choques monetários
inesperados jamais produzem efeitos reais sobre produto e emprego.
\end{enumerate}
Alternativas: i) apenas A e B; ii) apenas A, B e C; iii) apenas B e D;
iv) todas.

\begin{correctbox}[title={Resposta: ii) Apenas A, B e C --- com nota de ambiguidade na C}]
\textbf{A --- V.} Efeito Fisher: $\Delta M\uparrow \Rightarrow \pi^e\uparrow
\Rightarrow i\uparrow$ no longo prazo (um-para-um com preços flexíveis).

\textbf{B --- V.} Pela teoria das expectativas,
$i^n_t = \frac{1}{n}\sum_j \E_t\,i^1_{t+j} + \theta_{nt}$: anúncio crível de
política mais restritiva reduz a inflação esperada e as taxas curtas nominais
futuras esperadas --- as taxas \emph{longas} caem imediatamente, antes da
inflação efetiva ceder (espelho da conclusão dos slides: expansão derruba
curtas e sobe longas).

\textbf{C --- V.} Mesmo mecanismo da B: a queda das expectativas
inflacionárias reduz as taxas longas \emph{sem exigir movimento da Selic
corrente} (que é apenas um dos $n$ termos da média). Se B é verdadeira, C é
verdadeira pela mesma equação.

\textbf{D --- F.} O ``jamais'' inverte o resultado clássico: mesmo sob
expectativas racionais, choques monetários \emph{inesperados} têm efeitos
reais --- é justamente o componente não antecipado que produz efeito (Lucas),
e com rigidez nominal o efeito é ainda mais direto.

\emph{Nota de ambiguidade}: uma leitura literal de ``independentemente da taxa
Selic corrente'' na C (a taxa longa \emph{contém} a Selic corrente na média)
poderia levar o docente a considerá-la falsa --- caso em que a resposta seria
i) (apenas A e B). A leitura econômica padrão, porém, sustenta C verdadeira.
\end{correctbox}

% ==========================================================================
\newpage
\begin{center}
{\large\bfseries\color{orangerbc} Balanço da prova}\\[0.2em]
\rule{\linewidth}{0.6pt}
\end{center}

\begin{resultbox}[title={O que esta prova confirma sobre o padrão do docente}]
\begin{enumerate}[leftmargin=1.6em, itemsep=0.2em]
\item \textbf{Reciclagem intensa}: 7 das 12 questões reutilizam itens de
provas anteriores (Q2, Q4, Q6, Q7, Q8 com itens literais; Q3 e Q12
reformuladas);
\item \textbf{Inversão de uma palavra}: a Q3 trocou $\lambda=1$ por
$\lambda=0$ e acrescentou um ``não'' --- transformando o V de 2024 em F e o F
em V. Decorar \emph{letras} de gabarito não funciona; decorar o
\emph{conteúdo} sim;
\item \textbf{Absolutos derrubam itens}: ``necessariamente'' (Q7-ii),
``jamais'' (Q12-D), ``não sofre qualquer influência'' (Q7-iv) --- exceto o
caso ricardiano sobre T e D (Q2-iii, Q8-ii), em que a irrelevância \emph{é} o
teorema;
\item \textbf{Novidades vieram dos slides}: Q9 (síntese), Q10 (multiplicadores
negativos --- Aula 13), Q11 (Clarida --- novidade da Lista 3 2026), Q12
(estrutura a termo).
\end{enumerate}
\end{resultbox}

\begin{provabox}[title={Registro}]
Prova aplicada em 15/07/2026 (12 questões objetivas de escolha única, modelo
P1). Gabarito reconstruído de forma independente com base nas chaves oficiais
disponíveis; Q9 e Q12 sinalizadas com ambiguidade de redação. Conferência
detalhada da prova preenchida em
\texttt{lista\_III/auditoria\_nota\_p2\_2026.pdf}.
\end{provabox}

\vfill
{\small\color{grayrbc} Gabarito comentado elaborado para estudo pessoal ---
Macroeconomia I, PPGEco-UFES 2026. Não substitui o gabarito oficial.}

\end{document}
